\documentclass[aspectratio=1610]{beamer}%, handout
  \usepackage{graphics}
\usepackage{pifont}
\usepackage{ulem}
\usepackage{xcolor}
\usepackage{modifycolor}
\usepackage{lipsum}
\usepackage{fontspec}
\usepackage{caption}
\renewcommand{\figurename}{FIGURE}
\captionsetup[figure]{labelfont={color=leftFootlineColor, scriptsize}, textfont={color=normalBlockColor, scriptsize}}
\usepackage{tabularx}
\newcolumntype{Y}{>{\centering\arraybackslash}X}
\usepackage{tikz}
\usepackage{standalone}
\usepackage{svg}
\usepackage{multicol}
\usepackage{catchfilebetweentags}
\usepackage{xifthen}
\usetikzlibrary{arrows}
\usetikzlibrary{backgrounds}
\setmonofont[
  Contextuals={Alternate}
]{FuraCode Nerd Font}
\setsansfont{Roboto Medium}
\usepackage{pgfpages}


\usepackage[cache=false,outputdir=build]{minted}
\definecolor{bg_code}{HTML}{282828}
\usemintedstyle{darcula}
\setminted{
      fontsize=\scriptsize, 
      linenos,
      numbersep=0pt,
      gobble=5,
      framesep=3mm} 
      \renewcommand{\theFancyVerbLine}{\texttt{{{\arabic{FancyVerbLine}}}}}
\usepackage{adjustbox}
\usepackage{environ}%
\usepackage{tikz}
\usetikzlibrary{patterns}
\usepackage[french]{babel}
\usepackage{beamerthemesideblue}


\setbeamerfont{note page}{size=\footnotesize}
\setbeamertemplate{note page}[custom]
\setbeamercolor{note page}{bg=backgroundColor, fg=white}

\newcommand{\insertlicense}{\includegraphics[height=1cm]{img/by-sa.png}}
\title[]{La rétroingénierie appliquée à Android}
\subtitle{La traque aux traqueurs}
\author{Maxime Catrice}
\date{\today}
\titlegraphic{\includegraphics[height=0.95cm]{logo.png}\bigbreak\insertlicense}

\setbeamertemplate{title page}[default][colsep=-4bp,rounded=true,shadow=false]
\setbeamercolor{background canvas}{bg=backgroundColor}
\setbeamertemplate{blocks}[rounded][shadow=false]
\beamertemplatenavigationsymbolsempty
\newcounter{acolumn}%  Number of current column
\newlength{\acolumnmaxheight}%   Maximum column height
%%%%%%%%%%%%%%%%%%%%%%%%%


\makeatletter

% `column` replacement to measure height
\newenvironment{@acolumn}[1]{%
    \stepcounter{acolumn}%
    \begin{lrbox}{\@tempboxa}%
    \begin{minipage}{#1}%
}{%
    \end{minipage}
    \end{lrbox}
    \@tempdimc=\dimexpr\ht\@tempboxa+\dp\@tempboxa\relax
    % Save height of this column:
    \expandafter\xdef\csname acolumn@height@\roman{acolumn}\endcsname{\the\@tempdimc}%
    % Save maximum height
    \ifdim\@tempdimc>\acolumnmaxheight
        \global\acolumnmaxheight=\@tempdimc
    \fi
}

% `column` wrapper which sets the height beforehand
\newenvironment{@@acolumn}[1]{%
    \stepcounter{acolumn}%
    % The \autoheight macro contains a \vspace macro with the maximum height minus the natural column height
    \edef\autoheight{\noexpand\vspace*{\dimexpr\acolumnmaxheight-\csname acolumn@height@\roman{acolumn}\endcsname\relax}}%
    % Call original `column`:
    \orig@column{#1}%
}{%
    \endorig@column
}

% Save orignal `column` environment away
\let\orig@column\column
\let\endorig@column\endcolumn

% `columns` variant with automatic height adjustment
\NewEnviron{acolumns}[1][]{%
    % Init vars:
    \setcounter{acolumn}{0}%
    \setlength{\acolumnmaxheight}{0pt}%
    \def\autoheight{\vspace*{0pt}}%
    % Set `column` environment to special measuring environment
    \let\column\@acolumn
    \let\endcolumn\end@acolumn
    \BODY% measure heights
    % Reset counter for second processing round
    \setcounter{acolumn}{0}%
    % Set `column` environment to wrapper
    \let\column\@@acolumn
    \let\endcolumn\end@@acolumn
    % Finally process columns now for real
    \begin{columns}[#1]%
        \BODY
    \end{columns}%
}
\makeatother
%%%%%%%%%%%%%%%%%%%%%%%%%

\NewEnviron{frameNoSB}{
\makeatletter
\setbeamertemplate{sidebar canvas left}{}
\setbeamertemplate{sidebar left}{}
\makeatother
\begin{frame}
    \BODY
%\tableofcontents
\end{frame}
}

\NewEnviron{frameTitle}{
\setbeamertemplate{frametitle}[default][center]
\makeatletter
\setbeamertemplate{headline}{\color{backgroundColor}45\newline 45\newline 45\newline 45\newline 45\newline 45\newline 45\newline 45\newline 45\newline 45}
\setbeamertemplate{sidebar canvas left}{}
\setbeamertemplate{sidebar left}{}



\makeatother
\begin{frame}
\begin{minipage}[c]{\linewidth-\beamerleftmargin+\beamerrightmargin}
\BODY
\end{minipage}
\end{frame}
}

\setbeamerfont{title}{series=\bfseries,parent=structure}
\setbeamerfont{subtitle}{size=\Large,series=,parent=structure}

\makeatletter
\newlength\beamerleftmargin
\setlength\beamerleftmargin{\Gm@lmargin}

\newlength\beamerrightmargin
\setlength\beamerrightmargin{\Gm@rmargin}
\makeatother

\newcommand{\nologo}{\setbeamertemplate{logo}{}}

\newcommand{\slidetitle}[1][]{
  \frametitle{\insertsection\ifthenelse{\equal{#1}{}}{}{: #1}}
}
\newcommand{\mono}[1]{
\texttt{#1}
}

  \setbeameroption{show notes on second screen=top}
  \begin{document}

  \section{Légalité et rétroingénierie}
  \begin{frame}[t]
    \slidetitle[]%{Légalité et rétroingénierie}
    \note{\ExecuteMetaData[002-legal_notes.tex]{propIntelect}}
    \note{\onslide<5->{\ExecuteMetaData[002-legal_notes.tex]{loi}}}
    \note{\onslide<10->{\ExecuteMetaData[002-legal_notes.tex]{Resume}}}
%    \note{\onslide+<4->{\ExecuteMetaData[002-legal_notes.tex]{interropérabilité}}}
    \begin{columns}
      \begin{column}{0.7\linewidth}
        \begin{block}{Logiciels et propriété intellectuelle}<2->
          \begin{itemize}
          \item<3-> Logiciel protégeable
          \item<4-> Fonctionnalité en tant que telle non protégeable
          \end{itemize}
        \end{block}
        \begin{block}{Article 122-6-1 du code de la propriété intellectuelle}<5->
          \begin{itemize}
          \item<6-> Acquisition légale du logiciel
          \item<7-> Soit:
          \begin{itemize}
          \item<8-> La license ne l'interdit pas
          \item<9-> Réalisation à des fins d'interropérabilité
          \end{itemize} 
          \end{itemize}
        \end{block}
      \end{column}
      \begin{column}{0.25\linewidth}
        \includegraphics[width=0.95\linewidth]{img/justice.png}
      \end{column}
    \end{columns}
      \onslide+<10->{\begin{center}
        « On n'a donc pas le droit en France de démontrer techniquement qu'un logiciel présente des failles
  de sécurité, ou que la publicité pour ces logiciels est mensongère. Dormez tranquilles, citoyens, tous
  vos logiciels sont parfaits. » 
      \end{center}
      \vspace{-1.5\baselineskip}
      \hspace{0.66\linewidth}Guillermito
      }
    \end{frame}
  \end{document}

    \begin{frame}
    \slidetitle[Les Bug BugBounty]
    \note{\ExecuteMetaData[002-legal_notes.tex]{BugBounty}}
    \note{\onslide<4->{\ExecuteMetaData[002-legal_notes.tex]{Tweet}}}
    \pause
    \begin{block}{Qu'est ce qu'un Bug Bounty?}
      \centering
      \onslide<3->{Un bug bounty est un programme proposé par de nombreux sites web et développeurs de logiciel
       qui permet à des personnes de recevoir reconnaissance et compensation après avoir reporté
        des bugs, surtout ceux concernant des exploits et des vulnérabilités}
    \end{block}
    \vfill
    \centering
    \onslide<4->{
      \begin{figure}
        \includegraphics[width=0.5\linewidth]{tweet.png}
        \caption{Tweet montrant une faille dans macOS}
      \end{figure}}
      \end{frame}
  \end{document}
